\chapter{经典光场怎样在电子能量轴上写下相干边带}
\label{chap:feqo-classical-drive}

有限时间矩阵元能预测单次弱跃迁，但强一点的光场会让电子连续交换多个能量量子。若光场的量子涨落在当前精度下可忽略，可以先把它视为已知经典电场。电子沿轨迹积累相位；把这段相位按光学频率分解，就能同时得到整条边带，而不必逐阶求多光子微扰。

\section{沿轨迹的相位为何产生 Bessel 能量梳}
\label{chap:feqo-classical-pinem}

电子沿 \(+z\) 方向以中心速度 \(v_0\) 经过局域近场。光子诱导近场电子显微术（photon-induced near-field electron microscopy，PINEM）中，电子的离散能量边带正来自这样的相互作用。这里固定电场的记法：
\(E_z(z,t)=E_z^{(+)}(z)e^{-\ii\omega t}+E_z^{(+)*}(z)e^{\ii\omega t}\)。
若实验给出的复相量 \(\mathcal E_z\) 满足 \(E_z=\operatorname{Re}[\mathcal E_ze^{-\ii\omega t}]\)，则 \(E_z^{(+)}=\mathcal E_z/2\)。漏掉这个二分之一会把全部 Bessel 函数的自变量错放大一倍。

在 \(\phi=0\) 的规范中 \(E_z=-\partial_tA_z\)，所以正频矢势为 \(A_z^{(+)}=E_z^{(+)}/(\ii\omega)\)。要知道为何能沿一条轨迹积相位，先把第~\ref{chap:feqo-electron-state} 章的正能电子限制为沿 \(+z\) 运动的窄动量波包。令 \(\hat p_z=p_0+\delta\hat p\)，并把中心载波提出为 \(\Psi=e^{\ii(p_0z-E_0t)/\hbar}f(z,t)\)。自由色散给 \(E(p_0+\delta p)=E_0+v_0\delta p+O(\delta p^2/(\gamma_0^3m))\)；在正能分支投影后，Dirac 电流矩阵元 \(c\alpha_z\) 对窄包可用中心群速度 \(v_0\) 代替，故最小耦合项在这个阶次成为 \(-qv_0A_z\)。保留一阶动量偏差时，包络满足
\begin{equation}
\ii\hbar(\partial_t+v_0\partial_z)f(z,t)
=-qv_0A_z(z,t)f(z,t).
\label{eq:feqo-pinem-transport}
\end{equation}
沿特征线 \(z=v_0(t-t_0)\)，左边就是 \(\ii\hbar\dd f/\dd t\)，所以从入射到出射的解是作用势的指数。若包的动量宽度为 \(\delta p\)，这一步要求作用时间 \(T\) 内的自由二次色散相位 \(T(\delta p)^2/(2\gamma_0^3m\hbar)\ll1\)，并要求 \(|\delta p|/(\gamma_0^3mv_0)\ll1\) 使速度矩阵元近似一致。若实际占据到 \(|\ell|\leq N\)，也必须把 \(N\hbar|k_z|\) 计入动量偏离，并检查式~\eqref{eq:feqo-recoil-expansion} 的相位误差；若场强使 \(|qA_z|\) 接近电子动量尺度，或产生显著反射、横向偏转，还需回到原 Dirac 或完整波包方程。现在沿轨迹积分得到
\begin{align*}
U(t_0)&=\exp\!\left[-\frac{\ii}{\hbar}\int V(t)\dd t\right]
=\exp\!\left[\frac{\ii q}{\hbar}\int A_z(z,t_0+z/v_0)\dd z\right]\\
&=\exp\!\left[\beta e^{-\ii\omega t_0}-\beta^*e^{\ii\omega t_0}\right],
\end{align*}
其中无量纲复耦合定义为
\begin{equation}
\beta:=\frac{q}{\hbar\omega}\int_{-\infty}^{\infty}\dd z\,
E_z^{(+)}(z)e^{-\ii\omega z/v_0}.
\label{eq:feqo-beta}
\end{equation}
积分中的 \(qE_z^{(+)}\dd z\) 是能量；指数比较电子到达该位置时的光学相位。若场的空间因子是 \(e^{\ii kz}\) 且只在 \([-L/2,L/2]\) 非零，积分正比于 \(L\sinc[(k-\omega/v_0)L/2]\)：场越长，允许的波数失配越窄。

矩形窗的突然边界不是必需的。取一个可积分的局域近场 \(E_z^{(+)}(z)=E_*e^{-z^2/(2d^2)}e^{\ii kz}\)，其中 \(E_*\) 的单位是 \(\si{\volt\per\metre}\)，\(d\) 是衰减长度。设 \(\Delta k=k-\omega/v_0\)，完成平方并使用高斯积分，得到
\begin{equation}
\beta_G=\frac{qE_*\sqrt{2\pi}d}{\hbar\omega}
e^{-d^2\Delta k^2/2}.
\label{eq:feqo-gaussian-near-field-coupling}
\end{equation}
这一式同时提供量纲和极限检查：\(qE_*d/(\hbar\omega)\) 无量纲；\(\Delta k=0\) 时沿途场相位同向累积；\(|\Delta k|d\gg1\) 时正负贡献相消。可见同一积分耦合同时包含近场幅度、作用长度和速度失配；单一能谱不能把它们各自唯一分开。

再把两个相同的局域作用区放在 \(z=\pm D/2\)，令后一束光相对前一束多一个可控相位 \(\chi\)。变量替换 \(u=z-z_j\) 显示每区的积分多乘 \(e^{\ii\Delta k z_j}\)，因此在区间传播色散相位可忽略时
\[
\beta_{\rm two}=\beta_G
\left(e^{-\ii\Delta kD/2}+e^{\ii\chi}e^{\ii\Delta kD/2}\right),
\qquad
|\beta_{\rm two}|^2=4|\beta_G|^2
\cos^2\!\left(\frac{\Delta kD+\chi}{2}\right).
\]
调 \(\chi\) 可以使两次相位积累相长或相消，即使每一区单独都有非零耦合。若区间传播使不同电子能量分量取得不可忽略的相对相位，两个作用算符之间必须插入传播算符，上式的简单复数相加会失效。

写 \(\beta=|\beta|e^{\ii\phi}\)。由于指数是 \(2\ii|\beta|\sin(\phi-\omega t_0)\)，Jacobi--Anger 恒等式 \(e^{\ii x\sin u}=\sum_{\ell\in\mathbb Z}J_\ell(x)e^{\ii\ell u}\) 给
\begin{equation}
U(t_0)=\sum_{\ell\in\mathbb Z}
J_\ell(2|\beta|)e^{\ii\ell\phi}e^{-\ii\ell\omega t_0},
\qquad
c_\ell=J_\ell(2|\beta|)e^{\ii\ell\phi}.
\label{eq:feqo-pinem-sidebands}
\end{equation}
系数中 \(e^{-\ii\ell\omega t_0}\) 把中心能量 \(E_0\) 移到 \(E_0+\ell\hbar\omega\)；负 \(\ell\) 表示电子失去能量。第 \(\ell\) 条边带的概率是 \(P_\ell=J_\ell^2(2|\beta|)\)。它不含 \(\phi\)，但振幅含有；所以单次能谱无法读场相位，二次相干作用却可以。

回到刚算出的高斯近场，若 \(|\beta_G|\ll1\)，\(J_1(2|\beta_G|)=|\beta_G|+O(|\beta_G|^3)\)，故 \(P_{\pm1}=|\beta_G|^2+O(|\beta_G|^4)\)。双作用区的第一边带则是 \(4|\beta_G|^2\cos^2[(\Delta kD+\chi)/2]+O(|\beta_G|^4)\)。这条从场形状算到可见边带的链条也给出失效条件：\(|\beta_G|\) 不小时必须用完整 Bessel 函数，区间色散不可忽略时必须把传播插在两次作用中。

同一个 \(\beta\) 锁定每个 \(c_\ell\)；边带不能任意独立调节。若 \(2|\beta|\) 扫过 \(J_1\) 的一个零点，\(\ell=1\) 的人口会消失再出现，而总概率仍留在其他通道。这是多路径相干相消，不是光场停止与电子相互作用。

\begin{proposition}[Bessel 边带的归一化与两个矩]
对 \(P_\ell=J_\ell^2(2|\beta|)\)，有
\begin{equation}
\sum_\ell P_\ell=1,\qquad
\sum_\ell\ell P_\ell=0,\qquad
\sum_\ell\ell^2P_\ell=2|\beta|^2.
\label{eq:feqo-bessel-moments}
\end{equation}
\end{proposition}
\begin{proof}
令 \(u(\theta)=e^{\beta e^{-\ii\theta}-\beta^*e^{\ii\theta}}=\sum c_\ell e^{-\ii\ell\theta}\)。指数纯虚，故 \(|u|=1\)。周期 Fourier 正交性给 \(\sum|c_\ell|^2=(2\pi)^{-1}\int_0^{2\pi}|u|^2\dd\theta=1\)。由 \(J_{-\ell}(x)=(-1)^\ell J_\ell(x)\)，正负 \(\ell\) 的概率相等，故一阶矩为零。再将 \(\partial_\theta u=(-\ii\beta e^{-\ii\theta}-\ii\beta^*e^{\ii\theta})u\) 代入 Parseval 等式
\(\sum\ell^2|c_\ell|^2=(2\pi)^{-1}\int|\partial_\theta u|^2\dd\theta\)，交叉项的周期平均为零，余下 \(2|\beta|^2\)。
\end{proof}

这一结果要求电子在交互区主要前向传播，且实际占据边带 \(|\ell|\leq N\) 满足式~\eqref{eq:feqo-recoil-expansion} 下的相位误差远小于一。若横向偏转、反射或反冲可分辨，就需返回包含空间运动或精确反冲的方程；不能把 Bessel 公式解释为任意强场、任意几何的定律。

\begin{exercise}
从 \(J_0(x)=1-x^2/4+O(x^4)\)、\(J_1(x)=x/2+O(x^3)\) 求弱场下 \(P_0,P_{\pm1}\)，核对总概率到 \(|\beta|^2\) 阶为一。
\end{exercise}


单个电子的理想概率还不是实验中的柱状图。束斑内不同位置的电子、不同到达时刻的电子，会遇到不同的场强和相位。\(\beta\) 从一个固定数变成一个随机变量以后，必须先算条件概率，再对电子束平均。

\section{电子束为什么会洗浅单电子的振荡}
\label{chap:feqo-averaging}

式~\eqref{eq:feqo-pinem-sidebands} 是给定轨迹、给定到达时刻的条件概率。电子束内的每一个电子仍可经历幺正相位调制，却不一定经历同一个 \(\beta\)。设横向坐标 \(\mathbf R\) 与到达时刻 \(t_0\) 的联合概率密度 \(f(\mathbf R,t_0)\geq0\)，且 \(\int f\dd^2R\dd t_0=1\)。未逐个记录这两项时，计数比例为
\begin{equation}
\overline P_\ell=\int\dd^2R\dd t_0\,f(\mathbf R,t_0)
J_\ell^2\!\left(2|\beta(\mathbf R,t_0)|\right).
\label{eq:feqo-ensemble-pinem}
\end{equation}
这里先算概率再平均，因为不同电子是不同制备事件；\(J_\ell^2(2|\mathbb E\beta|)\) 通常不等于式~\eqref{eq:feqo-ensemble-pinem}。

一个可完整算出的例子能显示尺度。设电子束的横向分布为 \(f_R(R)=(2\pi\sigma_R^2)^{-1}e^{-R^2/(2\sigma_R^2)}\)，近场耦合为 \(\beta(R)=\beta_0e^{-R^2/(2L_E^2)}\)。在 \(|\beta_0|\ll1\) 时，\(J_1(2|\beta|)=|\beta|+O(|\beta|^3)\)，所以
\begin{align*}
\overline P_1
&=|\beta_0|^2\int_0^\infty\frac{R\dd R}{\sigma_R^2}
e^{-R^2/(2\sigma_R^2)-R^2/L_E^2}+O(|\beta_0|^4)\\
&=\frac{|\beta_0|^2}{1+2\sigma_R^2/L_E^2}+O(|\beta_0|^4).
\end{align*}
积分用了 \(\int_0^\infty Re^{-aR^2}\dd R=(2a)^{-1}\)。束斑宽度 \(\sigma_R\) 与场尺度 \(L_E\) 的比值直接控制可见边带；若 \(\sigma_R\gtrsim L_E\)，用中心电子的 \(\beta_0\) 拟合整个束流会系统性高估耦合均匀性。强场时仍须计算原始 Bessel 积分，不能延用这个弱场表达。

光脉冲与电子到达时间的平均是另一件事。若近场包络 \(\beta(t_0)=\beta_0e^{-t_0^2/(2\tau_L^2)}\)，到达时刻服从宽度 \(\sigma_t\) 的零均值高斯分布，同样在弱耦合下
\[
\overline P_1=\frac{|\beta_0|^2}{\sqrt{1+2\sigma_t^2/\tau_L^2}}+O(|\beta_0|^4).
\]
空间平均含二维 Jacobian，时间平均只含一维积分，两者的衰减因子不同。若还有探测器能量响应 \(R(y|E)\)，它在这些理想概率形成后再卷积；不能把所有对比度降低都叫作“能量展宽”。

强度平均与相位平均甚至改变不同的量。设每个电子都遇到同一个 \(|\beta|=b\)，但光学相位 \(\phi\) 随发射事件变化。由式~\eqref{eq:feqo-pinem-sidebands}，任一事件的能谱仍是 \(P_\ell=J_\ell^2(2b)\)，与 \(\phi\) 无关；因而均匀随机相位不会把这些峰直接洗浅。可是若把事件相位丢弃，边带密度矩阵的 \(\ell,\ell'\) 元会乘
\[
\langle e^{\ii(\ell-\ell')\phi}\rangle_\phi.
\]
例如相位在圆周上均匀分布时，该因子对 \(\ell\ne\ell'\) 为零；若相位抖动在一个参考相位附近服从方差 \(\sigma_\phi^2\) 的窄高斯分布，则模长是 \(e^{-(\ell-\ell')^2\sigma_\phi^2/2}\)。因此普通能谱可能几乎不变，后续时间聚束或参考场干涉却消失。要判断“对比度降低”究竟来自强度不均、相位不同步还是探测器分辨率，应比较不同类型的可观测量，不能仅凭一条能谱判断。

单一能谱只依赖 \(|\beta|\)。为了读出 \(\arg\beta\)，可让电子经过相位已知的第二个场。若两次相互作用之间的色散相位可忽略，乘积 \(U(\beta_2)U(\beta_1)=U(\beta_1+\beta_2)\)，其中指数在此平移不变模型里可交换。扫描 \(\beta_2=\beta_{\rm ref}e^{\ii\varphi}\) 后，\(|\beta_1+\beta_2|^2\) 含 \(2\operatorname{Re}(\beta_1\beta_2^*)\)，场相位遂转成人口变化。两区距离较大时必须把中间传播相位插入，简单相加不再成立。

\begin{exercise}
用二维高斯束的积分证明上式的 \(1+2\sigma_R^2/L_E^2\) 分母，并解释为什么在 \(\sigma_R\to0\) 时恢复单电子结果。
\end{exercise}


前两节把反冲相位视为可忽略；\cref{eq:feqo-recoil-expansion} 已说明该条件会随占据边带数增大而失败。保留精确通道能量后，哪些动量态能够同时参与将由失谐和耦合共同决定。

\section{有限反冲如何把多边格缩成少态动力学}
\label{chap:feqo-bragg}

低反冲模型中每条边在作用时间内都几乎同相。反冲相位可分辨后，\cref{eq:feqo-exact-recoil-chain} 的 \(\delta_\ell\) 把它们分开；只有接近共振的通道可以持续积累振幅。设目标态为 \(|0\rangle,|1\rangle\)，其余态组成 \(Q\) 子空间，目标投影为 \(P\)。这是正交投影：\(P^2=P=P^\dagger\)，\(Q=I-P\)。

本节的两种近似各占一块区域。\cref{fig:recoil-regime-map} 在跃迁面积 \(\mathscr A_{\rm hop}\) 与反冲相位 \(\omega_{\rm r}T\) 张成的对数平面上比较它们的总变差误差：左图的参照是理想 Bessel 梯，\(\omega_{\rm r}T\) 小时整条横带都接近色标下限；右图的参照是共振二态极限，误差小的一角只出现在 \(\mathscr A_{\rm hop}\lesssim0.3\)、\(\omega_{\rm r}T\gtrsim1\) 处，\(\epsilon=0.05\) 的等高线先在 \(\mathscr A_{\rm hop}\approx0.3\) 附近近乎竖直，再随 \(\omega_{\rm r}T\) 增大向右弯折。两个极限占据相反的区域，中间地带既不是 Bessel 也不是二态；右图左上方明暗相间的水平条纹还说明二态误差随 \(\omega_{\rm r}T\) 并不单调下降。

\begin{figure}[htbp]
\centering
\includegraphics[width=0.94\linewidth]{recoil_regime_map.pdf}
\caption{有限反冲格点上的近似精度分区。横纵坐标分别为跃迁面积 \(\mathscr A_{\rm hop}\) 与反冲相位 \(\omega_{\rm r}T\)，色标为总变差误差的 \(\log_{10}\epsilon\)。左图以理想 Bessel 梯 \(J_\ell^2(2\mathscr A_{\rm hop})\) 为参照，右图以共振二态极限为参照；两种比较都在 \(0\leftrightarrow1\) 精确共振的调谐下进行，等高线取 \(\epsilon=0.05\)，仅用于区分近似精度区域。}
\label{fig:recoil-regime-map}
\end{figure}

对不随时间变的旋转绘景哈密顿量 \(H\)，本征方程 \(H|\psi\rangle=z|\psi\rangle\) 的 \(Q\) 分量给 \((z-QHQ)Q|\psi\rangle=QHP P|\psi\rangle\)。只要 \(z\) 不在 \(QHQ\) 的谱上，便能解出 \(Q|\psi\rangle=(z-QHQ)^{-1}QHP P|\psi\rangle\)，代回 \(P\) 分量：
\begin{equation}
H_{\rm eff}(z)=PHP+PHQ(z-QHQ)^{-1}QHP.
\label{eq:feqo-feshbach-effective}
\end{equation}
这是消元恒等式，不要求弱耦合。它也说明被删掉的通道并非没有作用：第二项会移动能级并改变有效耦合。

三态例子可把量级算清。令目标两态能量为 \(0,\hbar\Delta\)，第三态能量为 \(\hbar\Delta_Q\)，且三态 哈密顿量 除以 \(\hbar\) 后为
\[
\begin{pmatrix}
0&g&h_0\\g^*&\Delta&h_1\\h_0^*&h_1^*&\Delta_Q
\end{pmatrix}.
\]
当 \(|\Delta_Q|\gg |g|,|h_0|,|h_1|,|\Delta|\) 时，在 \(z\) 接近零的范围展开 \((z-\hbar\Delta_Q)^{-1}=-(\hbar\Delta_Q)^{-1}+O(z/(\hbar^2\Delta_Q^2))\)，得到
\begin{equation}
\frac{H_{\rm eff}}{\hbar}=
\begin{pmatrix}0&g\\g^*&\Delta\end{pmatrix}
-\frac1{\Delta_Q}
\begin{pmatrix}|h_0|^2&h_0h_1^*\\h_1h_0^*&|h_1|^2\end{pmatrix}
+O\!\left(\frac{\|h\|^2\max(|g|,|\Delta|)}{|\Delta_Q|^2}\right).
\label{eq:feqo-bragg-two-level}
\end{equation}
这里 \(\|h\|^2=|h_0|^2+|h_1|^2\)。被消去态的振幅约为 \(\|h\|/|\Delta_Q|\)，人口约为这个比值的平方。若该比值不小，两态方程连人口都不能准确预测。二阶对角项若对两个目标态不同，还会移动共振位置，所以一般不能只凭裸失谐 \(\Delta=0\) 判共振。

谱选择之所以能发生，是因为共振调谐本身把外部格点推远了。\cref{fig:recoil-two-level-selection} 画出把 \(0\leftrightarrow1\) 调成简并后的失谐尺度：\(\ell=0,1\) 简并（图中标注 resonant），而二阶色散下最近的两个外部格点 \(\ell=-1,2\) 各自失谐 \(\delta_{-1}\approx\delta_2\approx2\omega_{\rm r}\)，所以它们对目标双态只有虚贡献。图中只画出谱位置这一层机制，定量的截断判断仍要用完整 \(Q\) 子空间的预解式。

\begin{figure}[htbp]
\centering
\includegraphics[width=0.62\linewidth]{recoil_two_level_selection.pdf}
\caption{调谐 \(0\leftrightarrow1\) 共振后的谱选择。四条横线给出二阶色散下 \(\ell=-1,0,1,2\) 的失谐尺度：中间两态简并（resonant），两侧的 \(\ell=-1,2\) 各失谐约 \(2\omega_{\rm r}\)。图只说明谱选择机制。}
\label{fig:recoil-two-level-selection}
\end{figure}

现在把上一章的反冲算例真正代入，而不是只说“保留两态”。令 \(\omega_d=k_zv_0+\omega_r\)、\(\Delta_Q=2\omega_r\)，并先取所有相邻边的耦合为同一个实数 \(g\)。按 \(|-1\rangle,|0\rangle,|1\rangle,|2\rangle\) 排列，保留四态时的矩阵为
\[
\frac H\hbar=
\begin{pmatrix}
\Delta_Q&g&0&0\\
g&0&g&0\\
0&g&0&g\\
0&0&g&\Delta_Q
\end{pmatrix}.
\]
取 \(P\) 为中间两态、\(Q\) 为外侧两态，便有 \(QHQ/\hbar=\Delta_QI\)、\(PHQ/\hbar=gI\)。代入\cref{eq:feqo-feshbach-effective}，在 \(|z|\ll\hbar\Delta_Q\) 处展开，得到
\[
\frac{H_{\rm eff}}\hbar=
\begin{pmatrix}0&g\\g&0\end{pmatrix}
-\frac{g^2}{\Delta_Q}I
+O\!\left(\frac{g^3}{\Delta_Q^2}\right).
\]
此处两个目标态的二阶位移恰好相同，只产生共同相位，\emph{不}移动它们的相对共振；前面的三态模型中的差分位移不是所有几何都必然出现。若两侧耦合分别为 \(g_-\) 与 \(g_+\)，二阶对角位移改为 \(-|g_-|^2/\Delta_Q\) 与 \(-|g_+|^2/\Delta_Q\)，其差才移动共振。这个对称性核对也提醒我们：消元不仅要估计项的阶数，还要保留实际耦合结构。

约化后的两态问题给出的正是失谐 Rabi 振荡，进入振荡频率的是含外部格点自能的 \(\Delta_{\rm eff}\) 而不是裸失谐 \(\Delta\)。\cref{fig:rabi-detuning} 把三条不同失谐的曲线放在一起：\(\Delta=0\) 时振幅为 1、首个极大在 \(|\Omega|t=\pi\)；\(\Delta/|\Omega|=1\) 时广义 Rabi 频率 \(\sqrt{|\Omega|^2+\Delta^2}\) 变大，极大提前到 \(|\Omega|t=\pi/\sqrt2\)，峰值按 \(|\Omega|^2/(|\Omega|^2+\Delta^2)=1/2\) 下降；\(\Delta/|\Omega|=2\) 时峰值进一步压到 \(1/5\)，且振荡在 \(|\Omega|t/\pi\) 尺度上明显变密。这正是后面用脉冲面积而不是时间标定 Bragg 转移的原因。

\begin{figure}[htbp]
\centering
\includegraphics[width=0.76\linewidth]{rabi_detuning.pdf}
\caption{谱隔离后的失谐 Rabi 振荡。横轴为以 \(\pi\) 为单位的演化时间 \(|\Omega|t/\pi\)，纵轴为激发态布居 \(P_{\rm e}\)，三条曲线分别取 \(\Delta/|\Omega|=0,1,2\)。失谐使振荡变快、峰值按 \(|\Omega|^2/(|\Omega|^2+\Delta^2)\) 下降。}
\label{fig:rabi-detuning}
\end{figure}

四态矩阵是为了让消元可手算，并非声称原无限梯只有四个态。若所有更远边也保留，它们要从目标态至少经过两次相邻跳跃才到达，所以在最近失谐有下界且 \(|g|/\omega_r\ll1\) 时，不改变刚才的二阶结果；接近其他高阶共振或耦合不小，则必须返回\cref{eq:feqo-exact-recoil-chain} 数值求解。

这里还要区分两类误差：截断只是不给全部格点，而格点自身的曲率会让无限梯上的整条分布一起变形。\cref{fig:pinem-recoil-breaking} 把曲率相位 \(\chi=\omega_{\rm curv}T\) 从零打开，取 \(\beta=2.5\)：\(\chi=0\) 时分布是理想 \(J_\ell^2(2|\beta|)=J_\ell^2(5)\)，在 \(\ell=\pm4\) 处约 0.153、\(\ell=\pm3\) 处约 0.133，而 \(\ell=\pm2\) 处几乎为零（约 0.002）；\(\chi=0.20\) 时 \(\ell=\pm4\) 与 \(\ell=\pm3\) 被压成约 0.148 的平肩，\(\ell=\pm2\) 的极小开始被填到 0.007；\(\chi=0.40\) 时最大值移到 \(\ell=\pm3\)（约 0.184，反而超过理想值），\(\ell=\pm4\) 降到 0.113，\(\ell=\pm5\) 从 0.068 塌到 0.020。峰位与峰高的这种重新分配说明，一旦曲率不可忽略，理想 Bessel 梯就不是“再多加几阶”能修正的对象。

\begin{figure}[htbp]
\centering
\includegraphics[width=0.76\linewidth]{pinem_recoil_breaking.pdf}
\caption{曲率相位 \(\chi=\omega_{\rm curv}T\) 对边带分布的影响（\(\beta=2.5\)）。实线为理想 Bessel 梯 \(J_\ell^2(2|\beta|)\)，另外两条曲线是加入二次曲率项 \(\chi\ell^2\) 后的精确幺正演化。\(\chi\) 增大时峰位内移、\(\ell=\pm2\) 处的极小被填高，理想梯的峰高比例不再成立。}
\label{fig:pinem-recoil-breaking}
\end{figure}

纵向光场交换以 \(p_z\to p_z+\hbar k_z\) 为主，反冲尺度为 \(\omega_r=\hbar k_z^2/(2\gamma_0^3m)\)。横向驻波若沿 \(x\) 周期，连接的是 \(p_x\to p_x+\hbar G\)，中心电子仍沿 \(z\) 前进；在 \(p_x\simeq0\) 处横向色散的二次系数是 \(1/(\gamma_0m)\)，对应 \(\omega_{r,\perp}=\hbar G^2/(2\gamma_0m)\)。二者的有效质量相差 \(\gamma_0^2\)，不可共用一个未说明方向的反冲频率。

若 \(N^2\omega_rT\ll1\)，实际占据的多条边不可分辨，出现薄光栅或 Raman--Nath 极限；若一条边满足 \(|\delta_{\ell+1}-\delta_\ell|T/\hbar\lesssim1\)，而所有其余边满足 \(|g|/|\Delta_{\rm gap}|\ll1\)，则\cref{eq:feqo-bragg-two-level} 描述 Bragg 少态区。两组条件都不满足时应数值求解原反冲格。

“耦合不小会怎样”可以直接测出来。取相邻格点耦合 \(g=\hbar\Omega/2\)，只保留两态的预言是 \(\sin^2(|\Omega|T/2)\)；\cref{fig:electron-bragg-convergence} 用完整反冲格的含时解（TDSE）检验它。左图扫过矩形脉冲的脉冲面积：\(r_\Omega=|\Omega|/\omega_{\rm r}=0.2\) 时曲线与二态虚线几乎重合，名义 \(\pi\) 脉冲处 \(P_1\approx0.995\)；\(r_\Omega=1\) 时峰值降到约 0.92 且峰位向小面积一侧移动；\(r_\Omega=5\) 时 \(P_1\) 最高只有约 0.35，要等到 \(|\Omega|T=2\pi\) 附近才回升到 0.75。右图把脉冲固定在名义 \(\pi\) 处扫描驱动比：\(r_\Omega\lesssim0.5\) 时 \(P_1\) 保持在 1，泄漏 \(P_{\rm leak}\) 从 \(r_\Omega\approx0.5\) 起明显上升，到 \(r_\Omega=7\) 时 \(P_1\) 只剩约 0.1 而 \(P_{\rm leak}\) 达到约 0.8。所以“只保留两个通道”与前面的 \(g/\omega_{\rm r}\ll1\) 是同一个条件。

\begin{figure}[htbp]
\centering
\includegraphics[width=0.94\linewidth]{electron_bragg_convergence.pdf}
\caption{矩形脉冲下反冲格转移的 TDSE 收敛性检验。左图取不同驱动比 \(r_\Omega=|\Omega|/\omega_{\rm r}\)，比较 \(0\to1\) 转移概率 \(P_1\) 随脉冲面积 \(|\Omega|T/\pi\) 的变化，虚线为零泄漏的两态预言 \(\sin^2(|\Omega|T/2)\)；右图在名义 \(\pi\) 脉冲处给出目标概率 \(P_1\) 与泄漏概率 \(P_{\rm leak}\) 随 \(r_\Omega\) 的变化。计算保留双侧格点 \(|\ell|\le25\)。}
\label{fig:electron-bragg-convergence}
\end{figure}

最后回到贯穿全节的二阶色散。\cref{fig:recoil-exact-vs-quadratic} 用 20 keV、800 nm 的参数，把相邻跃迁频率的精确相对论值 \([\mathcal E(p+(\ell+1)\hbar\kappa)-\mathcal E(p+\ell\hbar\kappa)]/\hbar\) 与二阶反冲近似 \(v\kappa+(2\ell+1)\hbar\kappa^2/(2\gamma_0^3m)\) 各自减去公共平移 \(v\kappa\) 后画在一起。在 \(|\ell|\le12\) 内两条曲线不可区分：\(\ell=-12\) 与 \(\ell=+12\) 处分别约为 \(-158\) 与 \(+171\) GHz，最大偏差不到 0.01 GHz，相对量级约 \(5\times10^{-5}\)。这一参数下 \(\hbar\kappa/p_0\approx4\times10^{-5}\)，所以二阶色散确实够用；该图同时也提醒，“二阶够用”是数值结论而非恒等式，\(\hbar\kappa/p_0\) 变大时必须重算。

\begin{figure}[htbp]
\centering
\includegraphics[width=0.76\linewidth]{recoil_exact_vs_quadratic.pdf}
\caption{相邻跃迁频率的精确相对论值与二阶反冲近似（20 keV 电子、800 nm 光）。横轴为格点编号 \(\ell\)，纵轴为扣去公共平移 \(v\kappa\) 后的频率差 \(\delta f_\ell\)（GHz）。两条曲线在 \(|\ell|\le12\) 内不可区分，最大偏差约 0.01 GHz。}
\label{fig:recoil-exact-vs-quadratic}
\end{figure}

\begin{exercise}
令 \(h_0=h_1=h\)、\(g=0\)，求消去第三态后产生的目标态间耦合及两个能级的共同移动，并检验 \(h/\Delta_Q\to0\) 极限。
\end{exercise}


当光场或结构周期延伸到许多周期时，还会出现空间平移和时间平移带来的标记简并。它们分别让动量和能量只在模一个周期量的意义下区分；为避免把记号上的简并误当成额外能量守恒，需分别说明这两种平移。

\section{空间周期与时间周期各自保存什么标记}
\label{chap:feqo-bloch-floquet}

若结构沿 \(z\) 每隔 \(a\) 重复，电子在 \(z\) 与 \(z+a\) 处看见相同势。定义平移算符 \((T_a\psi)(z)=\psi(z+a)\)。由 \(V(z+a)=V(z)\) 可直接核对 \([H,T_a]=0\)。先把体系放进长度为整数个 \(a\) 的大周期盒：此时能量子空间可按酉算符 \(T_a\) 的本征值分解，再取盒长趋于无穷的极限，所得无限结构的态可理解为广义本征态。因平移保持范数，\(T_a\) 的本征值模为一，写成 \(e^{\ii ka}\)。这就给出
\begin{equation}
\psi_k(z)=e^{\ii kz}u_k(z),\qquad u_k(z+a)=u_k(z).
\label{eq:feqo-bloch-state}
\end{equation}
证明只需令 \(u_k=e^{-\ii kz}\psi_k\)，再代入 \(\psi_k(z+a)=e^{\ii ka}\psi_k(z)\)。波数 \(k\) 与 \(k+G_m\)（\(G_m=2\pi m/a\)）对应相同平移本征值；周期函数相应乘 \(e^{-\ii G_mz}\)。这叫准动量，并不意味着电子真实动量只有一个值。

周期势的 Fourier 级数 \(V(z)=\sum_mV_me^{\ii G_mz}\) 立即显示矩阵元只连接 \(p\) 与 \(p+\hbar G_m\)。因此 Bragg 共振是这两个通道的自由能量近简并，再由 \(V_m\) 打开能隙；若 \(V_m=0\)，只有几何交点，不产生散射。

若 哈密顿量 随时间周期变化 \(H(t+T)=H(t)\)，普通定态 \(e^{-\ii Et/\hbar}|u\rangle\) 一般不再成立。定义一周期传播算符 \(U(T,0)\)。在有限维截断或它确有正规本征态的子空间内，取 \(U(T,0)|u_\alpha(0)\rangle=e^{-\ii\varepsilon_\alpha T/\hbar}|u_\alpha(0)\rangle\)；本征值模为一，因为传播算符酉。从该本征态出发，把精确解的相位 \(e^{-\ii\varepsilon_\alpha t/\hbar}\) 提出，余下部分每周期重复，得到
\begin{equation}
|\psi_\alpha(t)\rangle=e^{-\ii\varepsilon_\alpha t/\hbar}|u_\alpha(t)\rangle,
\qquad |u_\alpha(t+T)\rangle=|u_\alpha(t)\rangle.
\label{eq:feqo-floquet-state}
\end{equation}
\(\varepsilon\) 与 \(\varepsilon+m\hbar\Omega\) 给同一个一周期本征值，其中 \(\Omega=2\pi/T\)；这是准能的模 \(\hbar\Omega\) 自由度。式~\eqref{eq:feqo-floquet-state} 只要求周期性，不要求驱动很快。

上式的周期性可直接从传播律验证。周期驱动满足 \(U(t+T,T)=U(t,0)\)，所以 \(U(t+T,0)=U(t,0)U(T,0)\)。让它作用在选定的单周期本征态上，再乘 \(e^{\ii\varepsilon(t+T)/\hbar}\)，多出的 \(e^{-\ii\varepsilon T/\hbar}\) 与相位因子正好抵消，留下 \(|u(t+T)\rangle=|u(t)\rangle\)。若哈密顿量只有空间周期而没有时间周期，不能把这段证明借来谈准能；反过来时间周期也不会自动给电子真实动量一个倒格矢。

这里还有一个无限维的边界：酉算符可以有连续谱，未必有一组可数的正规本征向量张成整个空间。因此上式证明了每个存在的 Floquet 本征通道如何演化，并不宣称任意开放传播问题都能写成离散的 \(\sum_\alpha\)；连续谱部分应改用相应的谱积分。这个限定与 \(\varepsilon\) 模 \(\hbar\Omega\) 的记号自由度是两回事。

高频有效 哈密顿量 是额外近似。若非零 Fourier 分量 \(H_m\) 的相关矩阵元满足 \(\|H_m\|/(\hbar\Omega)\ll1\)，远离多光子共振，便可消去快速振荡到指定阶；被消去的周期内运动仍以微运动算符存在。若驱动频率接近两个能级差的整数分之一，即使 \(\Omega\) 数值很大，这个展开也可能因为小分母而失效。

\begin{exercise}
直接代入验证 \(k\to k+G_m\)、\(u_k\to e^{-\ii G_mz}u_k\) 不改变 Bloch 态；再验证 \(\varepsilon\to\varepsilon+m\hbar\Omega\) 时怎样重定义周期态。
\end{exercise}

